% Reused from chapters/09-vlan.tex
\begin{labvisual}{Recall the chapter model}
\begin{labgraphic}
\begin{tikzpicture}[node distance=0.9cm and 1.5cm]
\node[host] (pc1) {PC1};
\node[switch,right=1.5cm of pc1] (sw1) {SW1};
\node[switch,right=2.9cm of sw1] (sw2) {SW2};
\node[host,right=1.5cm of sw2] (pc3) {PC3};

\draw[flowarrow] (pc1) -- node[netlabel,above]{access} (sw1);
\draw[flowarrow] (sw1) -- node[netlabel,above]{trunk} (sw2);
\draw[flowarrow] (sw2) -- node[netlabel,above]{access} (pc3);

\node[frame,below=0.8cm of pc1,text width=1.9cm,align=center]
  {\textbf{no tag}\\ \tiny the host has no idea VLANs exist};
\node[frame,below=0.8cm of $(sw1)!0.5!(sw2)$,text width=2.5cm,align=center,
      draw=UTPBlue,fill=UTPPaleBlue]
  {\textbf{tag = VLAN 10}\\ \tiny added on entry, so the far switch knows};
\node[frame,below=0.8cm of pc3,text width=1.9cm,align=center]
  {\textbf{no tag}\\ \tiny stripped again on the way out};

\draw[->,draw=UTPMuted,dotted] ($(pc1)+(0,-0.35)$) -- ($(pc1)+(0,-0.75)$);
\draw[->,draw=UTPMuted,dotted] ($(sw1)!0.5!(sw2)+(0,-0.2)$) -- ($(sw1)!0.5!(sw2)+(0,-0.68)$);
\draw[->,draw=UTPMuted,dotted] ($(pc3)+(0,-0.35)$) -- ($(pc3)+(0,-0.75)$);
\end{tikzpicture}
\end{labgraphic}\end{labvisual}
